\documentclass[11pt]{article}
\usepackage[margin=1in]{geometry}
\usepackage{amsmath,amssymb,amsthm,mathrsfs}
\usepackage{graphicx}
\usepackage{hyperref}

\newtheorem{theorem}{Theorem}
\newtheorem{proposition}{Proposition}
\newtheorem{corollary}{Corollary}
\newtheorem{lemma}{Lemma}

\DeclareMathOperator{\dist}{dist}

\title{Nonlinear Disk Obstructions for the Bidisc Carath\'eodory-Fej\'er Problem}
\author{fa\_banach\_001 lane 5}
\date{2026-07-19}

\begin{document}
\maketitle

\section*{Status}

This is a partial result for arXiv:1508.07199, Rajeev Gupta,
\emph{The Carath\'eodory-Fej\'er Interpolation Problems and the von-Neumann
Inequality} \cite{Gupta1508}.  It gives additional necessary conditions for
the degree-two bidisc Carath\'eodory-Fej\'er problem and completely solves a
two-term subfamily.  It does not classify all pairs \(p_1,p_2\) in the CF
class.

\section*{Source Question}

In Section 4.4 of the source, after proving that the necessary condition
\(\|\mathscr T(M_{p_1},M_{p_2})\|\le 1\) is not sufficient in general, Gupta
asks what additional properties on \(p_1,p_2\) ensure membership in the CF
class.

\begin{center}
\includegraphics[width=0.95\linewidth]{figures/open_problem_crop.png}
\end{center}

For
\[
p(z_1,z_2)=a_{10}z_1+a_{01}z_2+a_{20}z_1^2+a_{11}z_1z_2+a_{02}z_2^2,
\]
the source uses
\[
p_1(\lambda)=a_{10}+a_{01}\lambda,\qquad
p_2(\lambda)=a_{20}+a_{11}\lambda+a_{02}\lambda^2.
\]
The source's explicit necessary condition is
\[
 |p_2(\lambda)|\le 1-|p_1(\lambda)|^2
 \qquad (|\lambda|=1).
\]

\section*{Proof Intuition}

The source's condition comes from testing an extension on the linear diagonal
disks \((z_1,z_2)=(t,\lambda t)\).  But a bidisc Schur function can also be
tested on nonlinear analytic disks.  The disks
\[
(z_1,z_2)=(\alpha t^2,\beta t)
\quad\text{and}\quad
(z_1,z_2)=(\alpha t,\beta t^2)
\]
mix one linear coefficient with one pure quadratic coefficient in the
second Taylor coefficient of a one-variable Schur function.  The tail \(q\)
of a CF extension has no total degree \(\le 2\), so it cannot change those
first two one-variable coefficients.  The elementary one-variable Schur
coefficient inequality therefore gives extra constraints.

\section*{One-Variable Input}

\begin{lemma}[Schur coefficient bound]\label{lem:schur}
Let \(F(t)=c_1t+c_2t^2+O(t^3)\) be holomorphic from \(\mathbb D\) to
\(\mathbb D\).  Then
\[
|c_2|\le 1-|c_1|^2.
\]
\end{lemma}

\begin{proof}
Since \(F(0)=0\), Schwarz's lemma gives \(F(t)=tG(t)\) with \(G\) in the
Schur class.  Then \(G(0)=c_1\) and \(G'(0)=c_2\).  The Schwarz-Pick lemma
applied to \(G\) at \(0\) gives \(|G'(0)|\le 1-|G(0)|^2\).
\end{proof}

\section*{Nonlinear Disk Obstructions}

\begin{theorem}\label{thm:nonlinear-obstructions}
Suppose there is a holomorphic \(q:\mathbb D^2\to\mathbb C\) such that
\[
q^{(k)}(0)=0\quad (|k|\le 2),
\qquad
\|p+q\|_{\mathbb D^2,\infty}\le 1.
\]
Then the coefficients of \(p\) satisfy
\[
|a_{10}|+|a_{02}|\le 1-|a_{01}|^2
\tag{1}
\]
and
\[
|a_{01}|+|a_{20}|\le 1-|a_{10}|^2.
\tag{2}
\]
\end{theorem}

\begin{proof}
Let \(f=p+q\).  Fix unimodular \(\alpha,\beta\) and set
\[
F_{\alpha,\beta}(t):=f(\alpha t^2,\beta t).
\]
This maps \(\mathbb D\) into \(\mathbb D\).  Since every monomial of \(q\) has
total degree at least \(3\), substituting \(z_1=\alpha t^2\), \(z_2=\beta t\)
produces only powers \(t^3,t^4,\ldots\) from \(q\).  Therefore
\[
F_{\alpha,\beta}(t)
=a_{01}\beta\,t+(a_{10}\alpha+a_{02}\beta^2)t^2+O(t^3).
\]
Lemma \ref{lem:schur} gives
\[
|a_{10}\alpha+a_{02}\beta^2|\le 1-|a_{01}|^2.
\]
Choosing \(\alpha\) and \(\beta\) so that the two summands have the same
argument gives (1).

The proof of (2) is symmetric.  For unimodular \(\alpha,\beta\), set
\[
G_{\alpha,\beta}(t):=f(\alpha t,\beta t^2).
\]
Again the tail \(q\) contributes no terms of order \(1\) or \(2\), and
\[
G_{\alpha,\beta}(t)
=a_{10}\alpha\,t+(a_{01}\beta+a_{20}\alpha^2)t^2+O(t^3).
\]
Lemma \ref{lem:schur} gives
\[
|a_{01}\beta+a_{20}\alpha^2|\le 1-|a_{10}|^2,
\]
and maximizing over phases gives (2).
\end{proof}

\section*{A Sharp Two-Term Classification}

\begin{corollary}\label{cor:two-term}
For \(a,b\in\mathbb C\), the polynomial
\[
p_{a,b}(z_1,z_2)=a z_1+b z_2^2
\]
has a degree-two bidisc CF extension if and only if
\[
|a|+|b|\le 1.
\]
\end{corollary}

\begin{proof}
Necessity is Theorem \ref{thm:nonlinear-obstructions} with
\[
a_{10}=a,\qquad a_{02}=b,\qquad a_{01}=a_{20}=a_{11}=0.
\]
For sufficiency, take \(q=0\).  If \(|a|+|b|\le 1\), then for
\((z_1,z_2)\in\mathbb D^2\),
\[
|a z_1+b z_2^2|
\le |a||z_1|+|b||z_2|^2
\le |a|+|b|
\le 1.
\]
Thus \(p_{a,b}\) itself is already a Schur-class extension.
\end{proof}

\begin{corollary}[Counterexample family to sufficiency of the source condition]
Let \(0<|a|<1\) and
\[
1-|a|<|b|\le 1-|a|^2.
\]
Then \(p_{a,b}(z_1,z_2)=a z_1+b z_2^2\) satisfies the source necessary
condition \(|p_2|+|p_1|^2\le 1\), but has no degree-two bidisc CF extension.
\end{corollary}

\begin{proof}
Here \(p_1(\lambda)=a\) and \(p_2(\lambda)=b\lambda^2\), so
\[
|p_2(\lambda)|+|p_1(\lambda)|^2=|b|+|a|^2\le 1.
\]
But \(|a|+|b|>1\), so Corollary \ref{cor:two-term} rules out a CF extension.
\end{proof}

\section*{Verification Notes}

The proof is analytic and uses only the Schwarz-Pick coefficient estimate.
The script \path{code/verify_disk_obstructions.py} performs a small numerical
sanity check: it samples random coefficients and phases to verify that the
maximized second coefficient on the nonlinear disks agrees with
\(|a_{10}|+|a_{02}|\) and \(|a_{01}|+|a_{20}|\), and it checks the two-term
classification boundary for sample values including the source example.
This code is not part of the proof.

The source PDF was checked at PDF page 73, printed page 61, where the example
\(p_1=1/\sqrt2\), \(p_2=z^2/2\) and the open question appear.

\section*{Limitations and Novelty Check}

This result is a partial answer to the source question.  It supplies additional
necessary conditions and a sharp classification of one two-term family, but it
does not give necessary and sufficient conditions for arbitrary quadratic
bidisc data.

Bounded novelty checks were performed in the local run indexes for
`1508.07199`, `CF class`, `completely polynomially extendible`, and the source
example \(p_1=1/\sqrt2,\ p_2=z^2/2\).  Web searches on 2026-07-19 found the
source thesis and the later Gupta-Misra Studia Mathematica paper, but no
separate record of this nonlinear-disk obstruction or the exact
\(a z_1+b z_2^2\) classification.  The novelty confidence is therefore
moderate, not definitive.

\section*{Human Review Recommendation}

Review Theorem \ref{thm:nonlinear-obstructions}.  The key point is that the
vanishing condition on \(q\) by total degree really prevents \(q\) from
altering the first two coefficients after the nonlinear disk substitutions.
Once that is checked, the rest is the one-variable Schwarz-Pick coefficient
bound and phase maximization.

\begin{thebibliography}{9}

\bibitem{Gupta1508}
R. Gupta,
\emph{The Carath\'eodory-Fej\'er Interpolation Problems and the von-Neumann
Inequality},
arXiv:1508.07199, 2015.

\bibitem{GuptaMisra2020}
R. Gupta and G. Misra,
\emph{The Carath\'eodory-Fej\'er interpolation on the polydisc},
Studia Math. 254 (2020), 265--293.

\end{thebibliography}

\end{document}
